\chapter{Metrics}

\section{Metrics for Vectors}

The metric of a vector \(a\) is described by different norms:

It also has the following properties:

\begin{itemize}

	\item \(\metric{a} \geq 0\) and \(\metric{a} = 0 \Leftrightarrow a = 0\)

	\item \(\metric{\lambda a} = |\lambda| \metric{a}\)

	\item \(\metric{a + b} \leq \metric{a} + \metric{b}\) (Triangle inequality)

	\item \(\metric{a + b}^2 = \metric{a}^2 + \metric{b}^2 + 2\langle a, b\rangle\) (Pythagorean theorem)

\end{itemize}

These are the most common metrics, although we will primarily be using the Euclidean metric:

\begin{itemize}

	\item \emph{Euclidean metric}:
	      \[
		      \metric{a} = \sqrt{\langle a, a\rangle} = \sqrt{a_1^2 + a_2^2 + \cdots + a_n^2}
	      \]
	\item \emph{Manhattan metric}:
	      \[
		      \metric{a}_1 = |a_1| + |a_2| + \cdots + |a_n|
	      \]
	\item \emph{Maximum metric}:
	      \[
		      \metric{a}_\infty = \max(|a_1|, |a_2|, \dots , |a_n|)
	      \]
\end{itemize}

\section{Metric for Matrices}

Given a metric \(\|\cdot\|\) over \(\R^n\), \(A \in \R^{n \times n}\), we define 

\[
    \metric{A} := \sup_{x \ne 0} \frac{\metric{Ax}}{\metric{x}} = \max_{\metric{x} = 1} \metric{Ax},
\]

as the \emph{matrix metric} for \(\R^{n \times n}\).

\begin{itemize}

    \item This norm is dependent on the chosen vector metric, hence it is called an \emph{induced matrix metric}.

    \item The vector metric and the induced metric norm are compatible, which means: 

    \[
        \metric{Ax}_{\text{vector}} \le \metric{A}_{\text{matrix}} \cdot \metric{x}_{\text{vector}}.
    \]

    \item For matrix metrics the following holds:

    \[
        \metric{AB} \le \metric{A} \cdot \metric{B}.
    \]

    \item For \(\metric{\R^{n \times n}}\), all matrix metrics are equivalent.

\end{itemize}

\subsection{Spectral Metric}

For \(\metric{\cdot}_2\) there is the so-called \emph{spectral metric}. 

\[
    \metric{A}_2 = \sqrt{\rho (A^T A)},
\]

where the \emph{spectral radius} \(\rho(B)\) is the eigenvalue of \(B\) with the greatest absolute value.

\subsection{Row Metric}

For \(\metric{\cdot}_\infty\) there is the \emph{row-sum metric} 

\[
    \metric{A}_\infty = \max_{i=1,\dots,n} \sum_{j = 1}^n |a_{ij}|.
\]

\subsection{Column Metric}

For \(\metric{\cdot}_1\) we get the \emph{column-sum metric} 

\[ 
    \metric{A}_1 = \max_{j=1,\dots,n} \sum_{i = 1}^n |a_{ij}|.
\]

\subsection{Frobenius Metric}

The \emph{Frobenius metric} given by 

\[
    \metric{A}_F = \sqrt{\dig(A^T A)}
\]

is not induced, but the following holds 

\[
    \metric{AB}_F \le \metric{A}_F \cdot \metric{B}_F,
\]

and, due to 

\[ 
    \metric{Ax}_2 \le \metric{A}_F \cdot \metric{x}_2, \forall x
\]

it is compatible with the vector metric \(\metric{\cdot}_2\). More precisely, \(\metric{A}_F\) is an upper bound for \(\metric{A}_2\).

